\documentclass[11pt,a4paper]{article}
% Shared preamble for RuPhO English problem statements (match T1 2026 style).
% Use: \input{latex/rupho_en_preamble.tex} after \documentclass.
\usepackage[margin=1in]{geometry}
\usepackage{amsmath,amssymb}
\usepackage{graphicx}
\usepackage{float}
\usepackage{enumitem}
\usepackage{microtype}
\usepackage{caption}
\usepackage{tabularx}
\usepackage{hyperref}

\captionsetup{font=small,labelfont=bf,skip=6pt}

\setlist[itemize]{leftmargin=1.2cm}
\setlist[enumerate]{leftmargin=1.2cm}

% One outer box per subproblem: [id | statement | points] (no inner rules)
\newcommand{\problembox}[3]{%
  \par\addvspace{0.42em}%
  \noindent
  \setlength{\fboxsep}{7.5pt}%
  \setlength{\fboxrule}{0.95pt}%
  \fbox{%
    \begin{minipage}{\dimexpr\linewidth-2\fboxsep-2\fboxrule\relax}
    \begin{tabularx}{\linewidth}{@{}>{\raggedright\arraybackslash\bfseries}p{2.1em} @{\hspace{0.9em}} >{\raggedright\arraybackslash}X @{\hspace{0.9em}} >{\raggedleft\arraybackslash\bfseries}p{2.4em}@{}}
    #1 & #2 & #3
    \end{tabularx}
    \end{minipage}%
  }%
  \par\addvspace{0.32em}%
}

\title{\textbf{Transient processes in a complex circuit}}
\author{\normalsize W22 (2022) --- English translation}
\date{}
\begin{document}
\maketitle

The electrical circuit shown in the figure consists of three capacitors $C_1=C_2=C_3=C$, two identical resistors $R_1=R_2=R$, a constant voltage source $\mathcal{E}$ with zero internal resistance and a switch $K$. While the key was open, the capacitors were not charged. Let us denote by $t$ the time after closing the key, the charges of the capacitors by $q_1$, $q_2$ and $q_3$ and the current strength in them ($I=dq/dt$) $I_1$, $I_2$ and $I_3$ in accordance with their numbers.

\begin{figure}[H]
    \centering
    \includegraphics[width=0.72\textwidth]{T2_EN_assets/fig1.jpg}
    \caption{Figure 1.}
\end{figure}

\problembox{A1}{Find the charges of capacitors $q_{1_0}$, $q_{2_0}$ and $q_{3_0}$ at $t\rightarrow0$.}{0.60}

\problembox{A2}{Find the charges of capacitors $q_{1_{\infty}}$, $q_{2_{\infty}}$ and $q_{3_{\infty}}$ at $t\rightarrow\infty$.}{0.60}

\problembox{A3}{Find the amount of heat $Q_0$ released in the system during time $t\rightarrow0$.}{0.40}

\problembox{A4}{Find the amount of heat $Q_{\infty}$ released in the system during time $t\rightarrow\infty$.}{0.40}

Let's move on to finding the time dependence of the charge of the first capacitor $q_1(t)$.

\problembox{A5}{From Kirchhoff's first law written for node $A$, express $I_2$ in terms of $I_1$, $q_1$, $R$ and $C$.}{0.50}

\problembox{A6}{From Kirchhoff's second law written for contour $ABCD$, express $q_2$ in terms of $I_1$, $q_1$, $R$ and $C$. Also find $I_2$ by differentiating $q_2$ with respect to time. Express your answer using $\cfrac{dI_1}{dt}$, $I_1$, $R$ and $C$.}{0.50}

\problembox{A7}{Equating the expressions for $I_2$ obtained in paragraphs A5 and A6, show that the differential equation describing the dependence $q_1(t)$ has the following form $$\textbackslash{}alpha\textbackslash{}frac{d^2q_1}{dt^2}+\textbackslash{}beta\textbackslash{}frac{dq_1}{dt}+q_1=0. $$   $\alpha$ and $\beta$  $R$ and $C$.}{0.40}

Further, consider it known that the solution to this differential equation must be sought in the form $$q_1(t)=Ae^{\textbackslash{}lambda{t}}, $$ where $\lambda$   from $R$ and $C$,  $A$ – from the initial conditions.

\problembox{A8}{Find $\lambda_1$ and $\lambda_2$ that satisfy the equation obtained in A7. Express your answers using $R$ and $C$. For $\lambda_1$, select the higher of the values.}{0.50}

In fact, the dependence $q_1(t)$ is represented as the sum of the found solutions $$q_1(t)=A_1e^{\textbackslash{}lambda_1{t}}+A_2e^{\textbackslash{}lambda_2{t}} $$

\problembox{A9}{From the initial conditions, find $A_1$ and $A_2$. Express your answers using $C$ and $\mathcal{E}$.}{0.80}

\problembox{A10}{Find the amounts of heat $Q_{R_1}$ and $Q_{R_2}$ released in the resistors during time $t\rightarrow\infty$. Express your answers using $C$ and $\mathcal{E}$.}{1.30}

$\textit{}$ : $\int e^{ax}dx=\frac{e^{ax}}{a}+C$

\vspace{1em}
\noindent\textit{Source:} \url{https://pho.rs/p/274}

\end{document}